\documentclass[11pt]{article}
\usepackage{amsmath, amssymb, geometry, amsthm, enumitem, xcolor}
\geometry{a4paper, margin=1in}

\title{\textbf{Homework 2.2 The Inverse of a Matrix}}
\author{Tashfeen Omran}
\date{September 2026}

\begin{document}

\maketitle

\tableofcontents
\newpage

\section{Part 1: Comprehensive Introduction, Context, and Prerequisites}

\subsection{Core Concepts}
In the realm of real numbers, every non-zero number $a$ has a multiplicative inverse, denoted as $a^{-1}$ or $\frac{1}{a}$, such that $a \cdot a^{-1} = 1$. The concept of \textbf{The Inverse of a Matrix} generalizes this foundational idea to linear algebra. 

For an $n \times n$ matrix $A$, an inverse is another $n \times n$ matrix, denoted $A^{-1}$, which satisfies the property:
\[ A A^{-1} = I \quad \text{and} \quad A^{-1} A = I \]
where $I$ is the $n \times n$ identity matrix (the matrix equivalent of the number 1, with 1s on the main diagonal and 0s elsewhere). 

If such a matrix $A^{-1}$ exists, $A$ is said to be \textbf{invertible} or \textbf{nonsingular}. If no such matrix exists, $A$ is called a \textbf{singular} matrix. It is crucial to note that only square matrices ($n \times n$) can have full inverses, and matrix division is undefined; instead of "dividing" by a matrix, we multiply by its inverse.

\subsection{Intuition and Derivation}
Why do we care about matrix inverses? Consider a standard algebraic equation: $5x = 10$. To isolate $x$, we multiply both sides by $5^{-1}$ (or divide by 5). 
Now consider a system of linear equations written in matrix form: $A\mathbf{x} = \mathbf{b}$. To solve for the vector $\mathbf{x}$, we cannot "divide" by $A$. However, if $A$ is invertible, we can multiply both sides on the left by $A^{-1}$:
\begin{align*}
    A^{-1}(A\mathbf{x}) &= A^{-1}\mathbf{b} \\
    (A^{-1}A)\mathbf{x} &= A^{-1}\mathbf{b} \quad \text{(by associativity)} \\
    I\mathbf{x} &= A^{-1}\mathbf{b} \\
    \mathbf{x} &= A^{-1}\mathbf{b}
\end{align*}
This provides an incredibly elegant, closed-form solution to systems of linear equations. 

For a $2 \times 2$ matrix, the inverse formula can be derived via algebraic elimination. By setting $A \cdot C = I$ and solving the resulting systems of equations for the components of $C$, we find that the inverse depends on a special scalar quantity known as the \textbf{determinant} ($ad - bc$).

\subsection{Historical Context and Motivation}
The formalization of matrix algebra, including the concept of the matrix inverse, largely occurred in the mid-19th century through the pioneering work of Arthur Cayley and James Joseph Sylvester. Mathematicians of the era were grappling with complex systems of linear equations arising in geometry and celestial mechanics. Cayley recognized that the coefficients of these systems could be detached and manipulated as single entities (matrices). The matrix inverse was developed to act as the ultimate "undo" button for linear transformations, streamlining the theoretical analysis of vector spaces and allowing for the generalization of algebraic rules to multi-dimensional spaces.

\subsection{Key Formulas}
\begin{itemize}
    \item \textbf{Definition of Inverse:} $A A^{-1} = I$ and $A^{-1} A = I$
    \item \textbf{Inverse of a $2 \times 2$ Matrix:} Let $A = \begin{bmatrix} a & b \\ c & d \end{bmatrix}$. If $ad - bc \neq 0$, then:
    \[ A^{-1} = \frac{1}{ad - bc} \begin{bmatrix} d & -b \\ -c & a \end{bmatrix} \]
    \item \textbf{Determinant of a $2 \times 2$ Matrix:} $\det(A) = ad - bc$
    \item \textbf{Inverse of a Product (Shoes and Socks Property):} $(AB)^{-1} = B^{-1} A^{-1}$
    \item \textbf{Inverse of a Transpose:} $(A^T)^{-1} = (A^{-1})^T$
    \item \textbf{Solving a System:} If $A\mathbf{x} = \mathbf{b}$ and $A$ is invertible, then $\mathbf{x} = A^{-1}\mathbf{b}$
\end{itemize}

\subsection{Prerequisites}
To master this topic, you must be completely fluent in:
\begin{itemize}
    \item \textbf{Matrix Multiplication:} Understanding how to take the dot product of rows and columns.
    \item \textbf{The Identity Matrix:} Knowing the structure and properties of $I$.
    \item \textbf{Row Reduction (Gauss-Jordan Elimination):} Essential for finding inverses of matrices larger than $2 \times 2$.
    \item \textbf{Matrix Matrix Equivalence:} Understanding that an invertible matrix $A$ is row equivalent to $I$.
\end{itemize}

\newpage
\section{Part 2: Detailed Homework Solutions}
\textit{Note: To provide the maximum depth of explanation while respecting formatting constraints, this section provides detailed, step-by-step solutions for the core foundational assignment: Computational Exercises 1-10, the True/False conceptual suite 11-18, and Exercise 39.}

\subsection{Exercise 1}
\textbf{Problem:} Find the inverse of the matrix $A = \begin{bmatrix} 8 & 3 \\ 5 & 2 \end{bmatrix}$.

\textbf{Solution:}
We use the $2 \times 2$ inverse formula. Let $a=8$, $b=3$, $c=5$, $d=2$.
First, calculate the determinant:
\[ \det(A) = ad - bc = (8)(2) - (3)(5) = 16 - 15 = 1 \]
Since $\det(A) \neq 0$, the inverse exists. Apply the formula:
\[ A^{-1} = \frac{1}{1} \begin{bmatrix} 2 & -3 \\ -5 & 8 \end{bmatrix} = \begin{bmatrix} 2 & -3 \\ -5 & 8 \end{bmatrix} \]
\textbf{Final Answer:} $A^{-1} = \begin{bmatrix} 2 & -3 \\ -5 & 8 \end{bmatrix}$

\subsection{Exercise 2}
\textbf{Problem:} Find the inverse of the matrix $A = \begin{bmatrix} 5 & 4 \\ 9 & 7 \end{bmatrix}$.

\textbf{Solution:}
Let $a=5$, $b=4$, $c=9$, $d=7$.
Calculate the determinant:
\[ \det(A) = (5)(7) - (4)(9) = 35 - 36 = -1 \]
Apply the formula:
\[ A^{-1} = \frac{1}{-1} \begin{bmatrix} 7 & -4 \\ -9 & 5 \end{bmatrix} = \begin{bmatrix} -7 & 4 \\ 9 & -5 \end{bmatrix} \]
\textbf{Final Answer:} $A^{-1} = \begin{bmatrix} -7 & 4 \\ 9 & -5 \end{bmatrix}$

\subsection{Exercise 3}
\textbf{Problem:} Find the inverse of the matrix $A = \begin{bmatrix} 8 & 3 \\ -7 & -3 \end{bmatrix}$.

\textbf{Solution:}
Let $a=8$, $b=3$, $c=-7$, $d=-3$.
Calculate the determinant:
\[ \det(A) = (8)(-3) - (3)(-7) = -24 - (-21) = -24 + 21 = -3 \]
Apply the formula:
\[ A^{-1} = \frac{1}{-3} \begin{bmatrix} -3 & -3 \\ 7 & 8 \end{bmatrix} = \begin{bmatrix} 1 & 1 \\ -7/3 & -8/3 \end{bmatrix} \]
\textbf{Final Answer:} $A^{-1} = \begin{bmatrix} 1 & 1 \\ -7/3 & -8/3 \end{bmatrix}$

\subsection{Exercise 4}
\textbf{Problem:} Find the inverse of the matrix $A = \begin{bmatrix} 3 & -2 \\ 7 & -4 \end{bmatrix}$.

\textbf{Solution:}
Let $a=3$, $b=-2$, $c=7$, $d=-4$.
Calculate the determinant:
\[ \det(A) = (3)(-4) - (-2)(7) = -12 - (-14) = -12 + 14 = 2 \]
Apply the formula:
\[ A^{-1} = \frac{1}{2} \begin{bmatrix} -4 & 2 \\ -7 & 3 \end{bmatrix} = \begin{bmatrix} -2 & 1 \\ -3.5 & 1.5 \end{bmatrix} \]
\textbf{Final Answer:} $A^{-1} = \begin{bmatrix} -2 & 1 \\ -7/2 & 3/2 \end{bmatrix}$

\subsection{Exercise 5}
\textbf{Problem:} Verify that the inverse you found in Exercise 1 is correct.

\textbf{Solution:}
To verify that $A^{-1}$ is correct, we must show that $A A^{-1} = I$.
\[ A = \begin{bmatrix} 8 & 3 \\ 5 & 2 \end{bmatrix}, \quad A^{-1} = \begin{bmatrix} 2 & -3 \\ -5 & 8 \end{bmatrix} \]
\[ A A^{-1} = \begin{bmatrix} 8 & 3 \\ 5 & 2 \end{bmatrix} \begin{bmatrix} 2 & -3 \\ -5 & 8 \end{bmatrix} \]
Row 1 $\cdot$ Col 1: $(8)(2) + (3)(-5) = 16 - 15 = 1$ \\
Row 1 $\cdot$ Col 2: $(8)(-3) + (3)(8) = -24 + 24 = 0$ \\
Row 2 $\cdot$ Col 1: $(5)(2) + (2)(-5) = 10 - 10 = 0$ \\
Row 2 $\cdot$ Col 2: $(5)(-3) + (2)(8) = -15 + 16 = 1$
\[ A A^{-1} = \begin{bmatrix} 1 & 0 \\ 0 & 1 \end{bmatrix} = I \]
\textbf{Final Answer:} The product $AA^{-1}$ yields the identity matrix, verifying the inverse is correct.

\subsection{Exercise 6}
\textbf{Problem:} Verify that the inverse you found in Exercise 2 is correct.

\textbf{Solution:}
We must show that $A A^{-1} = I$.
\[ A = \begin{bmatrix} 5 & 4 \\ 9 & 7 \end{bmatrix}, \quad A^{-1} = \begin{bmatrix} -7 & 4 \\ 9 & -5 \end{bmatrix} \]
\[ A A^{-1} = \begin{bmatrix} 5 & 4 \\ 9 & 7 \end{bmatrix} \begin{bmatrix} -7 & 4 \\ 9 & -5 \end{bmatrix} \]
Row 1 $\cdot$ Col 1: $(5)(-7) + (4)(9) = -35 + 36 = 1$ \\
Row 1 $\cdot$ Col 2: $(5)(4) + (4)(-5) = 20 - 20 = 0$ \\
Row 2 $\cdot$ Col 1: $(9)(-7) + (7)(9) = -63 + 63 = 0$ \\
Row 2 $\cdot$ Col 2: $(9)(4) + (7)(-5) = 36 - 35 = 1$
\[ A A^{-1} = \begin{bmatrix} 1 & 0 \\ 0 & 1 \end{bmatrix} = I \]
\textbf{Final Answer:} The product yields the identity matrix, verifying the inverse.

\subsection{Exercise 7}
\textbf{Problem:} Use the inverse found in Exercise 1 to solve the system:
\begin{align*}
    8x_1 + 3x_2 &= 2 \\
    5x_1 + 2x_2 &= -1
\end{align*}

\textbf{Solution:}
We can write this system in the matrix form $A\mathbf{x} = \mathbf{b}$:
\[ \begin{bmatrix} 8 & 3 \\ 5 & 2 \end{bmatrix} \begin{bmatrix} x_1 \\ x_2 \end{bmatrix} = \begin{bmatrix} 2 \\ -1 \end{bmatrix} \]
Since we know $A^{-1}$ from Exercise 1, we can multiply both sides on the left by $A^{-1}$:
\[ \mathbf{x} = A^{-1} \mathbf{b} \]
\[ \begin{bmatrix} x_1 \\ x_2 \end{bmatrix} = \begin{bmatrix} 2 & -3 \\ -5 & 8 \end{bmatrix} \begin{bmatrix} 2 \\ -1 \end{bmatrix} \]
Row 1 $\cdot$ Vector: $(2)(2) + (-3)(-1) = 4 + 3 = 7$ \\
Row 2 $\cdot$ Vector: $(-5)(2) + (8)(-1) = -10 - 8 = -18$
\[ \begin{bmatrix} x_1 \\ x_2 \end{bmatrix} = \begin{bmatrix} 7 \\ -18 \end{bmatrix} \]
\textbf{Final Answer:} $x_1 = 7, x_2 = -18$

\subsection{Exercise 8}
\textbf{Problem:} Use the inverse found in Exercise 2 to solve the system:
\begin{align*}
    5x_1 + 4x_2 &= -3 \\
    9x_1 + 7x_2 &= -5
\end{align*}

\textbf{Solution:}
The matrix form is $A\mathbf{x} = \mathbf{b}$:
\[ \begin{bmatrix} 5 & 4 \\ 9 & 7 \end{bmatrix} \begin{bmatrix} x_1 \\ x_2 \end{bmatrix} = \begin{bmatrix} -3 \\ -5 \end{bmatrix} \]
Using $A^{-1}$ from Exercise 2:
\[ \mathbf{x} = A^{-1} \mathbf{b} \]
\[ \begin{bmatrix} x_1 \\ x_2 \end{bmatrix} = \begin{bmatrix} -7 & 4 \\ 9 & -5 \end{bmatrix} \begin{bmatrix} -3 \\ -5 \end{bmatrix} \]
Row 1 $\cdot$ Vector: $(-7)(-3) + (4)(-5) = 21 - 20 = 1$ \\
Row 2 $\cdot$ Vector: $(9)(-3) + (-5)(-5) = -27 + 25 = -2$
\[ \begin{bmatrix} x_1 \\ x_2 \end{bmatrix} = \begin{bmatrix} 1 \\ -2 \end{bmatrix} \]
\textbf{Final Answer:} $x_1 = 1, x_2 = -2$

\subsection{Exercise 9}
\textbf{Problem:} Let $A = \begin{bmatrix} 1 & 2 \\ 5 & 12 \end{bmatrix}$, $\mathbf{b}_1 = \begin{bmatrix} -1 \\ 3 \end{bmatrix}$, $\mathbf{b}_2 = \begin{bmatrix} 1 \\ -5 \end{bmatrix}$, $\mathbf{b}_3 = \begin{bmatrix} 2 \\ 6 \end{bmatrix}$, $\mathbf{b}_4 = \begin{bmatrix} 3 \\ 5 \end{bmatrix}$. \\
a. Find $A^{-1}$, and use it to solve the four equations $A\mathbf{x} = \mathbf{b}_1$, $A\mathbf{x} = \mathbf{b}_2$, $A\mathbf{x} = \mathbf{b}_3$, $A\mathbf{x} = \mathbf{b}_4$. \\
b. The four equations in part (a) can be solved by the same set of row operations, since the coefficient matrix is the same in each case. Solve the four equations by row reducing the augmented matrix $[ A \quad \mathbf{b}_1 \quad \mathbf{b}_2 \quad \mathbf{b}_3 \quad \mathbf{b}_4 ]$.

\textbf{Solution (Part a):}
First, find $A^{-1}$. Let $a=1, b=2, c=5, d=12$.
$\det(A) = (1)(12) - (2)(5) = 12 - 10 = 2$.
\[ A^{-1} = \frac{1}{2} \begin{bmatrix} 12 & -2 \\ -5 & 1 \end{bmatrix} = \begin{bmatrix} 6 & -1 \\ -2.5 & 0.5 \end{bmatrix} \]
Now solve for each $\mathbf{x}_i = A^{-1} \mathbf{b}_i$:
\[ \mathbf{x}_1 = \begin{bmatrix} 6 & -1 \\ -2.5 & 0.5 \end{bmatrix} \begin{bmatrix} -1 \\ 3 \end{bmatrix} = \begin{bmatrix} -6 - 3 \\ 2.5 + 1.5 \end{bmatrix} = \begin{bmatrix} -9 \\ 4 \end{bmatrix} \]
\[ \mathbf{x}_2 = \begin{bmatrix} 6 & -1 \\ -2.5 & 0.5 \end{bmatrix} \begin{bmatrix} 1 \\ -5 \end{bmatrix} = \begin{bmatrix} 6 + 5 \\ -2.5 - 2.5 \end{bmatrix} = \begin{bmatrix} 11 \\ -5 \end{bmatrix} \]
\[ \mathbf{x}_3 = \begin{bmatrix} 6 & -1 \\ -2.5 & 0.5 \end{bmatrix} \begin{bmatrix} 2 \\ 6 \end{bmatrix} = \begin{bmatrix} 12 - 6 \\ -5 + 3 \end{bmatrix} = \begin{bmatrix} 6 \\ -2 \end{bmatrix} \]
\[ \mathbf{x}_4 = \begin{bmatrix} 6 & -1 \\ -2.5 & 0.5 \end{bmatrix} \begin{bmatrix} 3 \\ 5 \end{bmatrix} = \begin{bmatrix} 18 - 5 \\ -7.5 + 2.5 \end{bmatrix} = \begin{bmatrix} 13 \\ -5 \end{bmatrix} \]

\textbf{Solution (Part b):}
Set up the augmented matrix:
\[ \left[ \begin{array}{cc|cccc} 1 & 2 & -1 & 1 & 2 & 3 \\ 5 & 12 & 3 & -5 & 6 & 5 \end{array} \right] \]
Row replace: $R_2 \to R_2 - 5R_1$
\[ \left[ \begin{array}{cc|cccc} 1 & 2 & -1 & 1 & 2 & 3 \\ 0 & 2 & 8 & -10 & -4 & -10 \end{array} \right] \]
Scale: $R_2 \to \frac{1}{2}R_2$
\[ \left[ \begin{array}{cc|cccc} 1 & 2 & -1 & 1 & 2 & 3 \\ 0 & 1 & 4 & -5 & -2 & -5 \end{array} \right] \]
Row replace: $R_1 \to R_1 - 2R_2$
\[ \left[ \begin{array}{cc|cccc} 1 & 0 & -1 - 8 & 1 - (-10) & 2 - (-4) & 3 - (-10) \\ 0 & 1 & 4 & -5 & -2 & -5 \end{array} \right] \]
\[ \left[ \begin{array}{cc|cccc} 1 & 0 & -9 & 11 & 6 & 13 \\ 0 & 1 & 4 & -5 & -2 & -5 \end{array} \right] \]
The columns on the right match our $\mathbf{x}_i$ vectors from Part A. 

\textbf{Final Answer:} 
$\mathbf{x}_1 = \begin{bmatrix} -9 \\ 4 \end{bmatrix}$, $\mathbf{x}_2 = \begin{bmatrix} 11 \\ -5 \end{bmatrix}$, $\mathbf{x}_3 = \begin{bmatrix} 6 \\ -2 \end{bmatrix}$, $\mathbf{x}_4 = \begin{bmatrix} 13 \\ -5 \end{bmatrix}$.

\subsection{Exercise 10}
\textbf{Problem:} Use matrix algebra to show that if $A$ is invertible and $D$ satisfies $AD = I$, then $D = A^{-1}$.

\textbf{Solution:}
Given:
1. $A$ is invertible. By definition, this means there exists a unique matrix $A^{-1}$ such that $A^{-1}A = I$ and $AA^{-1} = I$.
2. $AD = I$.

We want to solve for $D$. Since $A^{-1}$ exists, we can left-multiply both sides of the equation $AD = I$ by $A^{-1}$:
\begin{align*}
    A^{-1}(AD) &= A^{-1}(I)
\end{align*}
Using the associative property of matrix multiplication:
\begin{align*}
    (A^{-1}A)D &= A^{-1}I
\end{align*}
Since $A^{-1}A = I$ and multiplying any matrix by the identity matrix leaves it unchanged ($A^{-1}I = A^{-1}$):
\begin{align*}
    I D &= A^{-1} \\
    D &= A^{-1}
\end{align*}
\textbf{Final Answer:} The derivation confirms $D$ must equal $A^{-1}$.

\subsection{Exercises 11-18 (True / False)}
\textbf{11. (T/F) In order for a matrix $B$ to be the inverse of $A$, both equations $AB = I$ and $BA = I$ must be true.}
\textbf{True.} The formal definition of an invertible matrix states there must be a matrix $C$ such that $AC = I$ and $CA = I$. (Note: for square matrices, one implies the other, but the definition requires both).

\textbf{12. (T/F) A product of invertible $n \times n$ matrices is invertible, and the inverse of the product is the product of their inverses in the same order.}
\textbf{False.} The inverse of a product is the product of their inverses in the \textit{reverse} order. Theorem 6(b): $(AB)^{-1} = B^{-1}A^{-1}$.

\textbf{13. (T/F) If $A$ and $B$ are $n \times n$ and invertible, then $A^{-1}B^{-1}$ is the inverse of $AB$.}
\textbf{False.} By Theorem 6(b), the inverse of $AB$ is $B^{-1}A^{-1}$. Matrix multiplication is not commutative.

\textbf{14. (T/F) If $A$ is invertible, then the inverse of $A^{-1}$ is $A$ itself.}
\textbf{True.} By Theorem 6(a), $(A^{-1})^{-1} = A$.

\textbf{15. (T/F) If $A = \begin{bmatrix} a & b \\ c & d \end{bmatrix}$ and $ad - bc \neq 0$, then $A$ is invertible.}
\textbf{True.} This is exactly Theorem 4. The determinant $ad - bc$ being non-zero is the condition for a $2 \times 2$ matrix to be invertible.

\textbf{16. (T/F) If $A = \begin{bmatrix} a & b \\ c & d \end{bmatrix}$ and $ad = bc$, then $A$ is not invertible.}
\textbf{True.} If $ad = bc$, then $ad - bc = 0$. By Theorem 4, a determinant of 0 means the matrix is singular (not invertible).

\textbf{17. (T/F) If $A$ is an invertible $n \times n$ matrix, then the equation $A\mathbf{x} = \mathbf{b}$ is consistent for each $\mathbf{b}$ in $\mathbb{R}^n$.}
\textbf{True.} By Theorem 5, the equation $A\mathbf{x} = \mathbf{b}$ has the unique solution $\mathbf{x} = A^{-1}\mathbf{b}$ for any $\mathbf{b}$, meaning it is always consistent.

\textbf{18. (T/F) If $A$ can be row reduced to the identity matrix, then $A$ must be invertible.}
\textbf{True.} Theorem 7 states that an $n \times n$ matrix $A$ is invertible if and only if it is row equivalent to the identity matrix $I_n$.

\subsection{Exercise 39}
\textbf{Problem:} Find the inverse of the matrix $\begin{bmatrix} 1 & 2 \\ 4 & 7 \end{bmatrix}$, if it exists.

\textbf{Solution:}
Let $A = \begin{bmatrix} 1 & 2 \\ 4 & 7 \end{bmatrix}$. 
$\det(A) = (1)(7) - (2)(4) = 7 - 8 = -1$.
Since $\det(A) \neq 0$, it is invertible.
\[ A^{-1} = \frac{1}{-1} \begin{bmatrix} 7 & -2 \\ -4 & 1 \end{bmatrix} = \begin{bmatrix} -7 & 2 \\ 4 & -1 \end{bmatrix} \]
\textbf{Final Answer:} $A^{-1} = \begin{bmatrix} -7 & 2 \\ 4 & -1 \end{bmatrix}$

\newpage
\section{Part 3: In-Depth Analysis of Problems and Techniques}

\subsection{A) Problem Types and General Approach}
The homework problems map directly onto four fundamental linear algebra archetypes:

\begin{enumerate}
    \item \textbf{Direct Formula Application (Problems 1-4, 39):} 
    \textit{Strategy:} These require immediate application of the $2 \times 2$ inverse formula. The strategy is straightforward: first, check if $\det(A) = ad - bc \neq 0$. If so, swap the main diagonal elements, flip the signs of the off-diagonal elements, and scalar-multiply by $\frac{1}{\det(A)}$.
    
    \item \textbf{Verification / Definition Testing (Problems 5-6):} 
    \textit{Strategy:} To prove a matrix is an inverse of another, you cannot just re-derive it. You must apply the definition: show that their product equals the Identity matrix ($I$). The strategy here is careful row-by-column dot product multiplication.
    
    \item \textbf{Solving Linear Systems (Problems 7-9):} 
    \textit{Strategy:} Translating a system of algebraic equations into the matrix form $A\mathbf{x} = \mathbf{b}$. Instead of row-reducing from scratch every time, if $A^{-1}$ is known, you instantly find the solution vector via $\mathbf{x} = A^{-1}\mathbf{b}$. Problem 9b introduced an advanced parallel strategy: row reducing multiple vectors simultaneously via $[ A \mid \mathbf{b}_1 \ \mathbf{b}_2 \ \mathbf{b}_3 \ \mathbf{b}_4 ]$.
    
    \item \textbf{Theoretical and Proof-Based (Problems 10, 11-18):} 
    \textit{Strategy:} These rely on manipulating symbols rather than numbers. The core technique is using valid matrix algebra rules (like left-multiplying both sides of an equation by the same matrix) and strictly applying theorems from the text (such as Theorem 6's reversal rule for inverse products).
\end{enumerate}

\subsection{B) Key Algebraic and Calculus Manipulations}
While linear algebra relies less on calculus techniques (like limits or integration) at this stage, the algebraic manipulations are highly specific and rigid:

\begin{itemize}
    \item \textbf{The Reversal Rule (Shoes and Socks Trick):} Seen in T/F 12 and 13. When you take the inverse of a product, you must reverse the order: $(AB)^{-1} = B^{-1}A^{-1}$. Think of it like getting dressed: you put your socks on ($B$), then your shoes ($A$). To reverse the process (take the inverse), you must first take off your shoes ($A^{-1}$), then your socks ($B^{-1}$).
    \item \textbf{Left/Right Multiplication Adherence:} Demonstrated in Problem 10. Because matrix multiplication is \textit{not} commutative ($AB \neq BA$), you must strictly respect sides. If you multiply the left side of an equation by $A^{-1}$ on the left, you \textit{must} multiply the right side by $A^{-1}$ on the left.
    \item \textbf{Block Matrix Row Reduction:} Demonstrated in Problem 9b. A massive simplification trick. Instead of solving $A\mathbf{x} = \mathbf{b}$ four separate times, you augment $A$ with a block matrix of all the $\mathbf{b}$ vectors. Row reducing the left side to $I$ simultaneously morphs the right side into the solution vectors. This is the exact same logic behind the algorithm for finding an inverse $[A \mid I] \sim [I \mid A^{-1}]$.
\end{itemize}

\newpage
\section{Part 4: "Cheatsheet" and Tips for Success}

\subsection{Crucial Formulas}
\begin{itemize}
    \item \textbf{$2 \times 2$ Inverse:} $A^{-1} = \frac{1}{ad - bc} \begin{bmatrix} d & -b \\ -c & a \end{bmatrix}$
    \item \textbf{General $n \times n$ Algorithm:} Row reduce $[ A \mid I ]$. If it reduces to $[ I \mid B ]$, then $B = A^{-1}$. If the left side cannot become $I$ (e.g., you get a row of all zeros), $A$ is singular.
    \item \textbf{Solving Systems:} $A\mathbf{x} = \mathbf{b} \implies \mathbf{x} = A^{-1}\mathbf{b}$ (valid only if $A$ is invertible).
    \item \textbf{Properties:} 
        \begin{itemize}
            \item $(A^{-1})^{-1} = A$
            \item $(AB)^{-1} = B^{-1}A^{-1}$
            \item $(A^T)^{-1} = (A^{-1})^T$
            \item $(cA)^{-1} = \frac{1}{c}A^{-1}$ (for a non-zero scalar $c$)
        \end{itemize}
\end{itemize}

\subsection{Shortcuts and Tricks}
\begin{itemize}
    \item \textbf{Quick Singularity Check:} For any $2 \times 2$ matrix, just cross-multiply the diagonals. If they equal each other ($ad = bc$), stop working immediately. The determinant is 0, and the inverse does not exist.
    \item \textbf{Fraction Deferral:} When calculating $2 \times 2$ inverses, leave the $\frac{1}{\det}$ scalar on the outside until the very end. If you distribute a fraction like $1/7$ into the matrix prematurely, all subsequent matrix multiplications will be miserable.
\end{itemize}

\subsection{Common Traps and Pitfalls}
\begin{itemize}
    \item \textbf{The Commutativity Trap:} Assuming $(AB)^{-1} = A^{-1}B^{-1}$. This is a fatal error in proofs. Always reverse the order!
    \item \textbf{The Division Trap:} Writing $\frac{B}{A}$ instead of $A^{-1}B$ or $BA^{-1}$. Matrix division does not exist. The side on which you multiply matters entirely.
    \item \textbf{The Addition Trap:} Assuming $(A + B)^{-1} = A^{-1} + B^{-1}$. This is profoundly false. (Think about it: $(1+1)^{-1} = 1/2$, but $1^{-1} + 1^{-1} = 2$).
\end{itemize}

\newpage
\section{Part 5: Conceptual Synthesis and The "Big Picture"}

\subsection{A) Thematic Connections}
The core theme of this topic is \textbf{Reversibility of Transformations}. Up to this point, a matrix $A$ was viewed statically as a grid of numbers or as a coefficient matrix. Now, we view $A$ as an active function (a linear transformation) mapping a vector $\mathbf{x}$ to a new vector $\mathbf{b}$. 

The inverse matrix $A^{-1}$ is the "undo" operator. If $A$ rotates a vector 90 degrees clockwise, $A^{-1}$ rotates it 90 degrees counter-clockwise. The mathematical search for inverses is thematically identical to finding inverse functions in pre-calculus (where $f^{-1}(f(x)) = x$) and anti-derivatives in calculus (where the integral reverses the derivative). 

\subsection{B) Forward and Backward Links}
\begin{itemize}
    \item \textbf{Backward Links:} The algorithm to find inverses relies entirely on Gauss-Jordan Elimination (Chapter 1). The idea that an invertible matrix must be row-equivalent to $I$ synthesizes everything you learned about pivot positions and basic/free variables.
    \item \textbf{Forward Links:} Matrix inverses are the bedrock for the \textbf{Invertible Matrix Theorem} (IMT), which you will see in Section 2.3. The IMT unites all concepts—determinants, linear independence, basis, span, and eigenvalues—into a single massive web of equivalencies. Furthermore, understanding inverses is required for Diagonalization and computing matrix exponentials later in the course.
\end{itemize}

\newpage
\section{Part 6: Real-World Application and Modeling}

\subsection{A) Concrete Scenarios (Finance, Stats, and Econ Focus)}
\begin{enumerate}
    \item \textbf{Modern Portfolio Theory (Finance):} In Harry Markowitz's Mean-Variance optimization, quantitative analysts try to construct a portfolio that minimizes risk (variance) for a given target return. The risk profile of the assets is stored in a massive Covariance Matrix ($\Sigma$). Finding the mathematically optimal portfolio weights requires inverting this covariance matrix ($\Sigma^{-1}$). 
    
    \item \textbf{Multiple Linear Regression (Statistics / Financial Econometrics):} When analysts build models to price securities (like the Arbitrage Pricing Theory or Fama-French models), they use historical data to find best-fit coefficients ($\beta$). The closed-form solution for the Ordinary Least Squares (OLS) estimator relies on taking the inverse of a matrix of data points.
    
    \item \textbf{Finite Difference Methods for Options Pricing (Quantitative Finance):} When solving the Black-Scholes Partial Differential Equation numerically (using an implicit method), the future value of the option is tied to the present value via a tridiagonal matrix $A$. Stepping backward in time to find the price today requires continuously computing $A^{-1}\mathbf{x}$.
\end{enumerate}

\subsection{B) Model Problem Setup: OLS Financial Modeling}
Imagine a quantitative analyst wants to model the return of a tech stock ($y$) based on three factors: the market return ($x_1$), interest rates ($x_2$), and inflation ($x_3$). They collect $N$ months of historical data.

\textbf{Variables:}
\begin{itemize}
    \item $\mathbf{y}$: An $N \times 1$ column vector of the tech stock's historical returns.
    \item $X$: An $N \times 4$ matrix where the first column is all 1s (for the intercept), and the subsequent columns are the historical data for $x_1, x_2, x_3$.
    \item $\hat{\beta}$: A $4 \times 1$ vector of the optimal coefficients we want to find.
\end{itemize}

\textbf{Mathematical Model:}
The goal of OLS is to minimize the sum of squared errors between the predicted returns and actual returns. Using multivariable calculus, taking the derivative with respect to the $\beta$ vector and setting it to zero yields the \textit{Normal Equations}:
\[ (X^T X) \hat{\beta} = X^T \mathbf{y} \]
To solve for the optimal weights $\hat{\beta}$, the analyst must rely on matrix inversion. Assuming the columns of $X$ are linearly independent, the matrix $(X^T X)$ is a square, invertible matrix. The computer must calculate:
\[ \hat{\beta} = (X^T X)^{-1} X^T \mathbf{y} \]

\newpage
\section{Part 7: Common Variations and Untested Concepts}

A concept highly related to inverses but not directly tested in this specific homework set is \textbf{The Adjugate Matrix (or Classical Adjoint)} for finding inverses of $3 \times 3$ matrices directly without row reduction.

\textbf{The Concept:}
While row reducing $[A \mid I]$ is algorithmically efficient for computers, there is a closed-form algebraic formula for the inverse of \textit{any} sized matrix using determinants. 
\[ A^{-1} = \frac{1}{\det(A)} \text{adj}(A) \]
Where $\text{adj}(A)$ is the transpose of the cofactor matrix. For a $3 \times 3$, this requires calculating 9 distinct $2 \times 2$ determinants!

\textbf{Brief Worked Example:}
Let $A = \begin{bmatrix} 1 & 0 & 0 \\ 0 & 2 & 0 \\ 0 & 0 & 3 \end{bmatrix}$. 
For diagonal matrices, the inverse is remarkably simple. We don't even need the full Adjugate. A standard variation problem asks you to invert a diagonal matrix.
Because $A \mathbf{x} = \mathbf{b}$ implies scaling each component directly, the inverse of a diagonal matrix is simply a diagonal matrix with the reciprocals of the entries on the diagonal:
\[ A^{-1} = \begin{bmatrix} 1/1 & 0 & 0 \\ 0 & 1/2 & 0 \\ 0 & 0 & 1/3 \end{bmatrix} \]
You can verify this visually: multiplying $A$ by $A^{-1}$ results in $\begin{bmatrix} 1 & 0 & 0 \\ 0 & 1 & 0 \\ 0 & 0 & 1 \end{bmatrix} = I$.

Another common variation is \textbf{Block Matrix Inversion}. Often in advanced statistics, we have a massive matrix broken into chunks: $M = \begin{bmatrix} A & 0 \\ 0 & B \end{bmatrix}$. Its inverse is elegantly $M^{-1} = \begin{bmatrix} A^{-1} & 0 \\ 0 & B^{-1} \end{bmatrix}$, provided $A$ and $B$ are invertible square blocks.

\newpage
\section{Part 8: Advanced Diagnostic Testing: "Find the Flaw"}

Below are 5 problems with subtle, fatal flaws. Your task as the grader is to spot the error, explain why it's mathematically invalid, and provide the correct solution.

\subsection{Problem 1: The Fast Formula Failure}
\textbf{Question:} Find the inverse of $A = \begin{bmatrix} 4 & 5 \\ 2 & 3 \end{bmatrix}$.
\\
\textbf{Flawed Solution:}
The determinant is $(4)(3) - (5)(2) = 12 - 10 = 2$.
To find the inverse, I flip the signs of the main diagonal and swap the off-diagonal elements.
\[ A^{-1} = \frac{1}{2} \begin{bmatrix} -4 & 2 \\ 5 & -3 \end{bmatrix} = \begin{bmatrix} -2 & 1 \\ 2.5 & -1.5 \end{bmatrix} \]
\\
\textbf{Flaw:} The student reversed the operations in the formula: they swapped the off-diagonals and negated the main diagonal, instead of swapping the main diagonal and negating the off-diagonals.
\\
\textbf{Correct Solution:}
\[ A^{-1} = \frac{1}{2} \begin{bmatrix} 3 & -5 \\ -2 & 4 \end{bmatrix} = \begin{bmatrix} 1.5 & -2.5 \\ -1 & 2 \end{bmatrix} \]

\subsection{Problem 2: The Order of Operations Error}
\textbf{Question:} If $A = \begin{bmatrix} 2 & 0 \\ 0 & 3 \end{bmatrix}$ and $B = \begin{bmatrix} 1 & 1 \\ 0 & 1 \end{bmatrix}$, find $(AB)^{-1}$.
\\
\textbf{Flawed Solution:}
By the properties of inverses, $(AB)^{-1} = A^{-1} B^{-1}$.
$A^{-1} = \begin{bmatrix} 1/2 & 0 \\ 0 & 1/3 \end{bmatrix}$ and $B^{-1} = \begin{bmatrix} 1 & -1 \\ 0 & 1 \end{bmatrix}$.
Therefore, $(AB)^{-1} = \begin{bmatrix} 1/2 & 0 \\ 0 & 1/3 \end{bmatrix} \begin{bmatrix} 1 & -1 \\ 0 & 1 \end{bmatrix} = \begin{bmatrix} 1/2 & -1/2 \\ 0 & 1/3 \end{bmatrix}$.
\\
\textbf{Flaw:} The student failed to reverse the order; the property is $(AB)^{-1} = B^{-1} A^{-1}$.
\\
\textbf{Correct Solution:}
$(AB)^{-1} = B^{-1} A^{-1} = \begin{bmatrix} 1 & -1 \\ 0 & 1 \end{bmatrix} \begin{bmatrix} 1/2 & 0 \\ 0 & 1/3 \end{bmatrix} = \begin{bmatrix} 1/2 & -1/3 \\ 0 & 1/3 \end{bmatrix}$.

\subsection{Problem 3: The Additive Illusion}
\textbf{Question:} Let $C = A + B$. Assume $A, B$, and $C$ are all invertible $n \times n$ matrices. Write an expression for $C^{-1}$.
\\
\textbf{Flawed Solution:}
Because matrix multiplication distributes over addition, the inverse distributes over addition as well.
$C^{-1} = (A + B)^{-1} = A^{-1} + B^{-1}$.
\\
\textbf{Flaw:} The inverse operation is highly non-linear and does absolutely not distribute over addition; $(A+B)^{-1} \neq A^{-1} + B^{-1}$.
\\
\textbf{Correct Solution:}
There is no simple, general expansion for $(A+B)^{-1}$. The expression remains $(A+B)^{-1}$ (unless one applies specialized formulas like the Woodbury matrix identity under specific rank updates).

\subsection{Problem 4: The Commutative Collision}
\textbf{Question:} Solve the matrix equation $XA = B$ for the unknown matrix $X$, assuming $A$ is invertible.
\\
\textbf{Flawed Solution:}
To isolate $X$, we must divide by $A$. This means multiplying both sides by $A^{-1}$.
$X = A^{-1}B$.
\\
\textbf{Flaw:} The student multiplied by $A^{-1}$ on the \textit{left} side of the right-hand equation, but to cancel the $A$ in $XA$, they needed to multiply on the \textit{right}.
\\
\textbf{Correct Solution:}
Multiply both sides on the right by $A^{-1}$:
$(XA)A^{-1} = BA^{-1} \implies X(AA^{-1}) = BA^{-1} \implies X = BA^{-1}$.

\subsection{Problem 5: The Premature Inverse}
\textbf{Question:} Find the inverse of $C = \begin{bmatrix} 1 & 2 & 3 \\ 4 & 5 & 6 \end{bmatrix}$.
\\
\textbf{Flawed Solution:}
We set up the augmented matrix $[C \mid I]$ and row reduce.
\[ \left[ \begin{array}{ccc|cc} 1 & 2 & 3 & 1 & 0 \\ 4 & 5 & 6 & 0 & 1 \end{array} \right] \]
Wait, the identity matrix must match the row dimension, so it's a $2 \times 2$ identity. Row reducing this will eventually yield a $2 \times 3$ matrix on the right, which is the inverse.
\\
\textbf{Flaw:} The student attempted to find the inverse of a non-square matrix.
\\
\textbf{Correct Solution:}
The matrix $C$ is $2 \times 3$. Only square ($n \times n$) matrices can be invertible. Therefore, $C$ has no inverse.

\end{document}